\section{Conclusion}
\label{sec:conclusion}

We presented \sentinel, a multi-layer defense architecture for LLM security
that synthesizes 58 security paradigms from 19 scientific domains into 7 novel
security primitives. Evaluated against 250,000 simulated attacks spanning 15
categories, the architecture achieves approximately 98.5\% detection and
containment---approaching the theoretical floor of 98--99\% bounded by proven
impossibility results. Rather than pursuing unattainable perfect security,
\sentinel~demonstrates that principled cross-domain synthesis can close the gap
between current ad hoc defenses and fundamental theoretical limits.

\paragraph{Key contributions.}
Our principal contributions are:

\begin{itemize}
  \item A systematic cross-domain paradigm synthesis methodology that analyzes
    58 paradigms drawn from 19 scientific domains, providing a repeatable
    framework for discovering security primitives at domain boundaries.

  \item Seven novel security primitives---\pasr, \cafl, \gps, \aas, \irm,
    \mire, and \tsa---of which five are genuinely new, confirmed via 51
    prior art searches returning zero existing implementations.

  \item A paradigm shift from backdoor detection---proven impossible by
    Goldwasser and Kim~\cite{goldwasser2022}---to architectural containment
    via \mire, rendering the impossibility result irrelevant by changing the
    defensive objective.

  \item \pasr---the first primitive that preserves information flow control
    through lossy semantic transformations, formalized as a provenance lifting
    functor (category-theoretic fibration) that maintains security labels
    across meaning-preserving rewrites.

  \item A demonstration that formal linguistics is the most underexploited
    domain for LLM security---an irony, given that large language models are,
    at their core, language models.
\end{itemize}

\paragraph{Broader impact.}
The cross-domain synthesis methodology itself may be the most significant
contribution of this work. By systematically searching across domain
boundaries, we discovered that mature solutions to the LLM security problem
already exist---scattered across unrelated fields, invisible to specialists
working within any single domain. Information flow control, adversarial
robustness, formal verification, and linguistic theory each hold pieces of the
puzzle; none alone is sufficient. We believe this methodology is applicable
beyond LLM security to other nascent fields where fundamental problems remain
unsolved, and where the barrier to progress is not missing science but missing
synthesis.

\paragraph{Availability.}
The Sentinel Core implementation (L1: 53 engines, 704 patterns) is open source
at \url{https://github.com/DmitrL-dev/AISecurity}. Implementation of the
remaining primitives is ongoing.
